% Block 2 of lista-reais-lectures-01-05 (input by the drivers in Listas/).
% Order and ratings: lista-reais-lectures-01-05/order.md.
\bloco{2}{Static games II: Nash equilibrium and mixed strategies}{Lecture 2}

% ---------------------------------------------------------------------------
\exercicio[1]{Hawk-Dove}{Foundations}{Tadelis (2013), \emph{Game Theory: An Introduction}, Ch.~5
``Pinning Down Beliefs: Nash Equilibrium'', Exercise 5.6 ``Hawk-Dove'', p.~96 [PDF 115]. No answer
in the book. Reproduced verbatim, with the book's footnote 12.}

\rot{5.6}\textbf{Hawk-Dove:} The following game has been widely used in evolutionary biology to
understand how fighting and display strategies by animals could coexist in a population. For a
typical Hawk-Dove game there are resources to be gained (e.g., food, mates, territories), denoted as
$v$. Each of two players can choose to be aggressive, as Hawk ($H$), or compromising, as Dove ($D$).
If both players choose $H$ then they split the resources but lose some payoff from injuries,
denoted as $k$. Assume that $k > \frac{v}{2}$. If both choose $D$ then they split the resources but
engage in some display of power that carries a display cost $d$, with $d < \frac{v}{2}$. Finally, if
player $i$ chooses $H$ while $j$ chooses $D$ then $i$ gets all the resources while $j$ leaves with no
benefits and no costs.
\begin{partes}
\item Describe this game in a matrix.
\item Assume that $v = 10$, $k = 6$, and $d = 4$. What outcomes can be supported as pure-strategy
Nash equilibria?$^{12}$
\end{partes}
{\small $^{12}$\,In the evolutionary biology literature, the analysis performed is of a very
different nature. Instead of considering the Nash equilibrium analysis of a static game, the
analysis is a dynamic one in which successful strategies ``replicate'' in a large population. This
analysis is part of a methodology called evolutionary game theory. For more on the subject see
Gintis (2000).\par}
\resolespaco

% ---------------------------------------------------------------------------
\exercicio[1]{Battle of the Sexes}{Foundations}{Muñoz-Garcia \& Toro-Gonzalez (2019), \emph{Strategy
and Game Theory: Practice Exercises with Answers}, 2nd ed., Ch.~2 ``Pure Strategy Nash
Equilibrium and Simultaneous-Move Games with Complete Information'', Exercise 2 ``Battle of the
Sexes''$^{\mathrm{A}}$, p.~41 [PDF 54]. Worked answer in the book, right after the statement.
Reproduced verbatim (statement only); Fig.~2.4, printed inside the answer on the same page,
rebuilt, with the book's row label ``Housband'' kept as printed.}

A husband and a wife are leaving work, and do not remember which event they are attending to
tonight. Both of them, however, remember that last night's argument was about either attending to
the football game (the most preferred event for the husband) or the opera (the most preferred event
for the wife). To make matters worse, their cell phones broke, so they cannot call each other to
confirm which event they are attending to. As a consequence, they must simultaneously and
independently decide whether to attend to the football game or the opera.

The payoff matrix in Fig.~2.4 describes the preference of the husband (wife) for the football game
(opera, respectively), but also indicates that both players prefer to be together rather than being
alone (even if they are alone at their most preferred event). Find the set of Nash Equilibria of
this game.
\jogo{Housband}{Wife}{\begin{tabular}{c|c|c|}
\hd{} & \hd{\emph{Football}} & \hd{\emph{Opera}}\\\cline{2-3}
\emph{Football} & 3, 1 & 0, 0\\\cline{2-3}
\emph{Opera} & 0, 0 & 1, 3\\\cline{2-3}
\end{tabular}}
\vspace{-6pt}\begin{center}\small\textit{Fig.~2.4} Battle of the sexes game (Normal-form
representation)\end{center}
\resolespaco

% ---------------------------------------------------------------------------
\exercicio[2]{Game of Chicken}{Foundations}{Muñoz-Garcia \& Toro-Gonzalez (2019), \emph{Strategy and
Game Theory: Practice Exercises with Answers}, 2nd ed., Ch.~3 ``Mixed Strategies, Strictly
Competitive Games, and Correlated Equilibria'', Exercise 1 ``Game of Chicken''$^{\mathrm{A}}$,
pp.~86--87 [PDF 99--100]. Worked answer in the book, right after the statement. Reproduced
verbatim (statement only).}

Consider the game of Chicken (depicted in Fig.~3.1), in which two players driving their cars
against each other must decide whether or not to swerve.

\textbf{Part (a)} Is there any strictly dominated strategy for Player 1? And for Player 2?

\textbf{Part (b)} What are the best responses for Player 1? And for Player 2?

\textbf{Part (c)} Can you find any pure strategy Nash Equilibrium (psNE) in this game?

\textbf{Part (d)} Find the mixed strategy Nash Equilibrium (msNE) of the game. \emph{Hint}: denote by
$p$ the probability that Player 1 chooses Straight and by $(1 - p)$ the probability that he chooses
to Swerve. Similarly, let $q$ denote the probability that Player 2 chooses Straight and $(1 - q)$
the probability that she chooses to Swerve.

\textbf{Part (e)} Show your result of part (d) by graphically representing every player $i$'s best
response function $BRF_i(s_j)$, where $s_j = \{Swerve, Straight\}$ is the strategy selected by player
$j \neq i$.
\jogo{\emph{Player 1}}{\emph{Player 2}}{\begin{tabular}{c|c|c|}
\hd{} & \hd{Straight} & \hd{Swerve}\\\cline{2-3}
Straight & 0,0 & 3,1\\\cline{2-3}
Swerve & 1,3 & 2,2\\\cline{2-3}
\end{tabular}}
\vspace{-6pt}\begin{center}\small\textit{Fig.~3.1} Game of Chicken\end{center}
\resolespaco

% ---------------------------------------------------------------------------
\exercicio[2]{A mixed equilibrium by indifference}{Exam level}{Gibbons (1992), \emph{Game Theory for
Applied Economists}, Ch.~1 ``Static Games of Complete Information'', \S1.5 Problems (Section 1.3),
Problem 1.12, p.~51 [PDF 59]. No answer in the book. Reproduced verbatim.}

\rot{1.12.}Find the mixed-strategy Nash equilibrium of the following normal-form game.
\begin{center}\begin{tabular}{c|c|c|}
\hd{} & \hd{$L$} & \hd{$R$}\\\cline{2-3}
$T$ & 2, 1 & 0, 2\\\cline{2-3}
$B$ & 1, 2 & 3, 0\\\cline{2-3}
\end{tabular}\end{center}
\resolespaco

% ---------------------------------------------------------------------------
\exercicio[3]{The $n$-Firm Cournot Model}{Exam level}{Tadelis (2013), \emph{Game Theory: An
Introduction}, Ch.~5 ``Pinning Down Beliefs: Nash Equilibrium'', Exercise 5.8 ``The $n$-Firm
Cournot Model'', pp.~96--97 [PDF 115--116]. No answer in the book. Reproduced verbatim (the printed
``$Q = q_i + \cdots + q_n$'' is kept as printed).}

\rot{5.8}\textbf{The $n$-Firm Cournot Model:} Suppose there are $n$ firms in the Cournot oligopoly
model. Let $q_i$ denote the quantity produced by firm $i$, and let $Q = q_i + \cdots + q_n$ denote
the aggregate production. Let $P(Q)$ denote the market clearing price (when demand equals $Q$) and
assume that the inverse demand function is given by $P(Q) = a - Q$, where $Q \le a$. Assume that
firms have no fixed cost and that the cost of producing quantity $q_i$ is $cq_i$ (all firms have the
same marginal cost, and assume that $c < a$).
\begin{partes}
\item Model this as a normal-form game.
\item What is the Nash (Cournot) equilibrium of the game in which firms choose their quantities
simultaneously?
\item What happens to the equilibrium price as $n$ approaches infinity? Is this familiar?
\end{partes}
\resolespaco

% ---------------------------------------------------------------------------
\exercicio[3]{Cournot with asymmetric costs}{Exam level}{Gibbons (1992), \emph{Game Theory for Applied
Economists}, Ch.~1 ``Static Games of Complete Information'', \S1.5 Problems (Section 1.2), Problem
1.6, p.~49 [PDF 57]. No answer in the book. Reproduced verbatim.}

\rot{1.6.}Consider the Cournot duopoly model where inverse demand is $P(Q) = a - Q$ but firms have
asymmetric marginal costs: $c_1$ for firm 1 and $c_2$ for firm 2. What is the Nash equilibrium if
$0 < c_i < a/2$ for each firm? What if $c_1 < c_2 < a$ but $2c_2 > a + c_1$?

\begin{referencia}{``the Cournot duopoly model'': Gibbons (1992), \S1.2.A ``Cournot Model of
Duopoly'', p.~15 [PDF 23], the paragraph that sets up the model. Reproduced verbatim; footnote~6
(a pointer to later sections) omitted.}
Let $q_1$ and $q_2$ denote the quantities (of a homogeneous product) produced by firms 1 and 2,
respectively. Let $P(Q) = a - Q$ be the market-clearing price when the aggregate quantity on the
market is $Q = q_1 + q_2$. (More precisely, $P(Q) = a - Q$ for $Q < a$, and $P(Q) = 0$ for
$Q \ge a$.) Assume that the total cost to firm $i$ of producing quantity $q_i$ is $C_i(q_i) = cq_i$.
That is, there are no fixed costs and the marginal cost is constant at $c$, where we assume
$c < a$. Following Cournot, suppose that the firms choose their quantities simultaneously.
\end{referencia}
\resolespaco

% ---------------------------------------------------------------------------
\exercicio[3]{Homogeneous Bertrand}{Exam level}{Gibbons (1992), \emph{Game Theory for Applied
Economists}, Ch.~1 ``Static Games of Complete Information'', \S1.5 Problems (Section 1.2), Problem
1.7, pp.~49--50 [PDF 57--58]. No answer in the book. Reproduced verbatim. The pointer to Section
1.2.B (Bertrand with differentiated products) is context; the statement defines the whole
homogeneous-products model.}

\rot{1.7.}In Section 1.2.B, we analyzed the Bertrand duopoly model with differentiated products. The
case of homogeneous products yields a stark conclusion. Suppose that the quantity that consumers
demand from firm $i$ is $a - p_i$ when $p_i < p_j$, 0 when $p_i > p_j$, and $(a - p_i)/2$ when
$p_i = p_j$. Suppose also that there are no fixed costs and that marginal costs are constant at
$c$, where $c < a$. Show that if the firms choose prices simultaneously, then the unique Nash
equilibrium is that both firms charge the price $c$.
\resolespaco

% ---------------------------------------------------------------------------
\exercicio[4]{Market Entry}{Exam level}{Tadelis (2013), \emph{Game Theory: An Introduction}, Ch.~6
``Mixed Strategies'', Exercise 6.8 ``Market Entry'', p.~125 [PDF 144]. No answer in the book.
Reproduced verbatim.}

\rot{6.8}\textbf{Market Entry:} Three firms are considering entering a new market. The payoff for
each firm that enters is $\frac{150}{n}$, where $n$ is the number of firms that enter. The cost of
entering is 62.
\begin{partes}
\item Find all the pure-strategy Nash equilibria.
\item Find the symmetric mixed-strategy equilibrium in which all three players enter with the same
probability.
\end{partes}
\resolespaco

% ---------------------------------------------------------------------------
\exercicio[4]{A unique IESDS survivor is a Nash equilibrium}{Advanced}{Mas-Colell, Whinston \& Green
(1995), \emph{Microeconomic Theory}, Ch.~8 ``Simultaneous-Move Games'', Exercises, Exercise
8.D.2$^{\mathrm{B}}$, p.~263 [PDF 142, right half]. Answer in the MWG solutions manual
(\texttt{Livros/}). Reproduced verbatim.}

\rot{8.D.2$^{\mathrm{B}}$}Show that if there is a unique profile of strategies that survives iterated
removal of strictly dominated strategies, this profile is a Nash equilibrium.
\resolespaco

% ---------------------------------------------------------------------------
\exercicio[5]{Existence in $2\times2$ games}{Advanced}{Mas-Colell, Whinston \& Green (1995),
\emph{Microeconomic Theory}, Ch.~8 ``Simultaneous-Move Games'', Exercises, Exercise
8.D.6$^{\mathrm{B}}$, p.~264 [PDF 143, left half]. Answer in the MWG solutions manual
(\texttt{Livros/}; 8.D.6 is there). Reproduced verbatim.}

\rot{8.D.6$^{\mathrm{B}}$}Consider any two-player game of the following form (where letters
indicate arbitrary payoffs):
\jogo{Player 1}{Player 2}{\begin{tabular}{c|c|c|}
\hd{} & \hd{$b_1$} & \hd{$b_2$}\\\cline{2-3}
$a_1$ & $u, v$ & $\ell, m$\\\cline{2-3}
$a_2$ & $w, x$ & $y, z$\\\cline{2-3}
\end{tabular}}
Show that a mixed strategy Nash equilibrium always exists in this game. [\emph{Hint:} Define player
1's strategy to be his probability of choosing action $a_1$ and player 2's to be his probability of
choosing $b_1$; then examine the best-response correspondences of the two players.]
\resolespaco
